% ============================================================================
% Chapter 4: PokeAgent Baseline Harness
% ============================================================================

\newpage
\section{PokeAgent Baseline Harness}\label{sec:pokeagent}

This chapter traces the development of the \pokeagent{} harness from its simplest competition-era form through to the expert system we deploy for the AutoEvolve experiments in Chapter~\ref{sec:autoevolve}. The narrative follows a chronological arc across three eras: (1)~the \emph{competition era}, in which a minimal four-module agent participated in the NeurIPS 2025 challenge; (2)~the \emph{benchmark-paper era}, in which an intermediate harness (internally called \emph{Evelord}) introduced structured objectives, a knowledge base, and A* pathfinding to produce the baseline results reported in the \pokeagent{} Challenge Retrospective~\citep{karten2025pokeagent}; and (3)~the \emph{expert harness} $\mathcal{H}_{\text{expert}}$, a full multi-agent orchestration system engineered for the long-horizon setting and designed to serve as the human-crafted ceiling against which AutoEvolve is measured.

We begin by establishing the baselines---human and model---against which all harness configurations are evaluated. We then describe the competition-era agent, review the early performance results that motivated each round of development, and conclude with a detailed account of the expert harness and the abstractions that distinguish it.

Figure~\ref{fig:ch4_neurips_leaderboard} presents the final standings of the NeurIPS 2025 speedrunning competition. The winning team, Heatz, used Scripted Policy Distillation (SPD): an LLM decomposes the task into subgoals and generates scripted policies for each, which are then distilled into neural networks via imitation learning and refined with reinforcement learning. The resulting agent completed the route in 40:13---nearly twice as fast as the second-place Hamburg Pok\'{e}Runners, who used a pure RL approach based on recurrent PPO with milestone-conditioned rewards. The remaining completions all used LLM harness architectures with varying levels of tool integration.

A key cross-track insight from the competition is that specialist methods---RL and search---consistently outperform generalist LLM approaches. In battling, the top participants all used RL or Monte Carlo tree search rather than LLM reasoning. In speedrunning, the top two finishers used RL-based methods, with Heatz's 40:13 more than $2\times$ faster than the best pure LLM harness approach. However, LLMs played a crucial enabling role in the winning approach: the LLM provided the prior (task structure and initial behavior) while RL provided refinement (faster execution and strategy discovery). This pattern---using LLMs for high-level reasoning and RL for low-level optimization---appears broadly applicable beyond games.

\begin{figure}[h]
\centering
\includegraphics[width=\linewidth]{ch3/track2_progress_time.pdf}
\caption{\textbf{NeurIPS Speedrunning Leaderboard.} Milestone progress vs.\ wall-clock time (left) and cumulative agent steps (right) for the six teams that completed all 15 milestones. Unlike most game benchmarks that pause between actions, our environment runs in real time---the game world continues while the agent reasons, making inference latency a first-class cost. This distinction reveals that Deepest (Judge's Choice) completed in the fewest steps (649) despite ranking 5th by time, highlighting the tradeoff between inference speed and action efficiency.}
\label{fig:ch4_neurips_leaderboard}
\end{figure}

% ---------------------------------------------------------------------------
\subsection{Speedrunning Baselines}\label{sec:baselines}

Before analyzing the evolution of the \pokeagent{} harness, we establish the reference points against which all systems are measured.

\subsubsection{Human Baselines}\label{sec:human_baselines}

Human speedrunners provide the upper and lower bounds for agent comparison in the tutorial setting. Our top human speedrunner completed the route from Littleroot Town to defeating Roxanne in 18 minutes, representing the upper bound on what is achievable with perfect game knowledge and motor execution. At the other end, the average human player---defined as a non-speedrunner familiar with Pok\'{e}mon games but not optimizing for speed---required approximately 1 hour 22 minutes (1:22:05). These two reference points frame the competitive landscape: the best frontier models complete this segment in under three hours, placing them roughly $2\times$ slower than the average human and an order of magnitude slower than expert human speedrunners. Closing this gap requires improvements in both decision quality and computational throughput.

\subsubsection{Harness vs.\ Model Capability}\label{sec:harness_vs_model}

A central challenge in evaluating LLM-based game agents is attribution: does observed performance stem from the underlying foundation model or from the surrounding harness? Prior efforts---Claude Plays Pok\'{e}mon~\citep{anthropic2025visible}, Gemini Plays Pok\'{e}mon~\citep{zhang2025geminiplayspokemon}, GPT Plays Pok\'{e}mon~\citep{engadget2025gptplayspokemon}---conflated these factors, making cross-system comparison scientifically problematic. To disentangle them, the \pokeagent{} Challenge introduced a five-dimensional characterization framework $(S, T, M, F, \Phi)$ that decomposes any agent system along orthogonal axes~\citep{karten2025pokeagent}:

\begin{itemize}[nosep]
    \item \textbf{State Representation $(S)$}: what the agent observes (visual frames, parsed RAM, structured text).
    \item \textbf{Tools $(T)$}: what actions the agent can take beyond raw button presses (pathfinding, knowledge retrieval, code execution).
    \item \textbf{Memory $(M)$}: how the agent maintains context across steps (raw context window, file-based summaries, persistent stores).
    \item \textbf{Feedback $(F)$}: whether the agent can evaluate its own decisions (pure observe-act vs.\ multi-agent self-reflection).
    \item \textbf{Fine-tuning $(\Phi)$}: whether model weights are updated (zero-shot, few-shot, RL fine-tuning).
\end{itemize}

\begin{figure}[h]
\centering
\includegraphics[width=\linewidth]{ch4/system_comparison.png}
\caption{\textbf{RPG System Comparison.} Analysis of Pok\'{e}mon RPG AI systems using the $(S, T, M, F, \Phi)$ framework, characterized at the version each system was initially popularized. Heterogeneity across dimensions makes direct performance comparison methodologically unsound without standardization.}
\label{fig:ch4_system_comparison}
\end{figure}

Figure~\ref{fig:ch4_system_comparison} applies this framework to contemporary Pok\'{e}mon RPG systems. The heterogeneity is striking: state representations range from raw visual frames to full RAM extraction; tool availability ranges from basic button inputs to full planning suites with code execution; memory architectures span raw context windows to engineered summarization pipelines. Even the reported evaluation metrics---variously ``steps,'' ``actions,'' and ``turns''---lack shared definitions across projects. The \pokeagent{} Speedrunning Track addresses precisely this gap by standardizing the state interface, tool access, and evaluation protocol, enabling principled harness $\times$ model comparisons. We use this framework throughout the remainder of this chapter and in Chapter~\ref{sec:autoevolve} to precisely characterize each harness configuration.

% ---------------------------------------------------------------------------
\subsection{Early Development: Competition Era}\label{sec:early_development}

The earliest \pokeagent{} harness---the version deployed during the NeurIPS 2025 competition---was a minimal system built to validate whether frontier VLMs could make meaningful progress in Pok\'{e}mon Emerald at all. The codebase consisted of roughly 120 files, with the agent logic concentrated in six Python modules under a single \texttt{agent/} directory. In the $(S, T, M, F, \Phi)$ framework, this system occupied the simplest end of each axis: visual frames plus structured state text for $S$; no external tools for $T$ (the agent could only output raw GBA button presses); a rolling short-term buffer for $M$; a simple observe-plan-act loop for $F$; and zero-shot inference for $\Phi$.

\subsubsection{Four-Module Pipeline}\label{sec:four_module}

The competition-era agent implemented a linear pipeline of four modules, each backed by a separate VLM call:

\begin{enumerate}[nosep]
    \item \textbf{Perception}: receives the $240 \times 160$ game frame and structured state text; produces a scene analysis classifying the current context (overworld, battle, dialog, menu) and a flag indicating whether the situation requires deliberate reasoning.
    \item \textbf{Planning}: receives the memory context and current plan; checks whether the existing plan has been accomplished and, if so, generates a new strategic plan covering immediate goals, short-term objectives, and long-term strategy.
    \item \textbf{Memory}: maintains a rolling buffer of 50 entries (observation--action pairs) via a Python \texttt{deque}. Extracts significant events (map transitions, battle entries/exits) and formats the recent history into a context string. This module involves no LLM calls---it is purely programmatic.
    \item \textbf{Action}: receives all accumulated context (memory, plan, observation, frame) and produces a list of up to 10 GBA button presses (\texttt{A}, \texttt{B}, \texttt{UP}, \texttt{DOWN}, \texttt{LEFT}, \texttt{RIGHT}, \texttt{START}, \texttt{SELECT}).
\end{enumerate}

Each step through this pipeline required 3--4 VLM calls, making the agent slow on a per-action basis. Navigation was performed entirely by the LLM interpreting a text-based ASCII map grid and a movement preview indicating what lay in each cardinal direction---there was no algorithmic pathfinding.

\subsubsection{SimpleAgent Variant}\label{sec:simple_agent}

In practice, a second architecture within the same codebase became the primary workhorse: the \texttt{SimpleAgent}, a 1,315-line class that collapsed the four-module pipeline into a single VLM call per step. The prompt was structured to elicit chain-of-thought reasoning across five sections: \textsc{Analysis} $\to$ \textsc{Objectives} $\to$ \textsc{Plan} $\to$ \textsc{Reasoning} $\to$ \textsc{Action}. The model's response was parsed to extract each section, with the action section yielding button presses.

\begin{figure}[t]
\centering
\includegraphics[width=0.85\linewidth]{ch4/simple_agent_draft.png}
\caption{\textbf{PokeAgent Simple: Competition-Era Architecture.} The four-module pipeline (Perception $\to$ Planning $\to$ Memory $\to$ Action) processes a game frame and structured state through sequential VLM calls, producing button presses. The SimpleAgent variant collapses this into a single structured VLM call. Neither variant has access to external tools or persistent memory. \tersoo{Refine this draft figure with a cleaner diagram.}}
\label{fig:ch4_simple_agent}
\end{figure}

The SimpleAgent expanded the memory buffer to 100 state entries and 50 recent actions (of which 30 and 40, respectively, were surfaced to the LLM), added basic stuck detection (flagging repeated actions at the same coordinates after 8 consecutive steps), and introduced a rudimentary objective system: pre-loaded storyline milestones supplemented by dynamic sub-objectives that the LLM could create or complete via inline commands (\texttt{ADD\_OBJECTIVE:}, \texttt{COMPLETE\_OBJECTIVE:}). The system prompt was remarkably minimal---just two sentences describing the agent's role---with all substantive prompt engineering embedded in the per-module prompt templates.

This architecture had no MCP tools, no pathfinding algorithm, no persistent knowledge base, and no mechanism for self-reflection. The agent navigated purely by interpreting the ASCII map, relying on the foundation model's world knowledge and spatial reasoning to infer appropriate movement. Despite these limitations, it was sufficient to complete the tutorial setting, establishing a proof of concept that frontier VLMs could play Pok\'{e}mon Emerald with minimal scaffolding.

% ---------------------------------------------------------------------------
\subsection{Early Performance}\label{sec:early_performance}

\subsubsection{Competition Era Results}\label{sec:competition_results}

The NeurIPS speedrunning competition (Figure~\ref{fig:ch4_neurips_leaderboard}) produced six teams that achieved 100\% milestone completion, with completion times ranging from 40:13 (Heatz) to 2:46:48. The diversity of approaches is instructive. At the top, Heatz's Scripted Policy Distillation used an LLM to decompose the task into subgoals with executable success conditions, generate scripted policies for each, and distill these into neural networks refined via RL---achieving emergent behaviors including BFS-based local pathfinding and trainer battle skipping~\citep{karten2025pokeagent}. Hamburg Pok\'{e}Runners took second with a pure RL approach: recurrent PPO with a CNN-MLP encoder, LSTM for temporal reasoning, and a binary milestone vector that served as both memory and goal conditioning. Anthonys (third place) combined deterministic A* pathfinding with phase-based prompt templates, decomposing the game into seven distinct phases with specialized navigation strategies. Deepest received the Judge's Choice award for sample efficiency, completing the route in just 649 actions (fewest of any team) using a training-free architecture built around a guidebook system, image-based goal recognition, and adaptive reasoning budgets. Full methodology descriptions for each team are provided in Appendix~\ref{sec:appendix_competitor_methods}.

The PokeAgent Simple baseline---the four-module / SimpleAgent architecture described in Section~\ref{sec:early_development}---was the weakest of the LLM harness approaches. Several observations from this era spurred subsequent development. First, the release of Gemini 3 models led to a notable increase in performance, particularly around visual processing: actions appeared more visually grounded, with less overreliance on structured text state~\citep{google2025gemini3}. Tasks requiring fine motor control, such as setting the clock in the player's bedroom, became more consistent. This underscored that base model capability---especially multimodal reasoning---was a first-order driver of performance. Second, the four-module architecture was poor on both step and time efficiency, requiring 3--4 VLM calls per action. Third, the agent lacked coherence across extended play: it would frequently become confused about game progress and infer its next action solely from the current screen and its world knowledge, making slow but eventual progress through the tutorial's largely linear structure.

\subsubsection{Benchmark Paper Results}\label{sec:benchmark_results}

Following the competition, the \pokeagent{} harness underwent its first major revision. The intermediate version---referred to internally as \emph{Evelord}---introduced three architectural changes that substantially improved consistency:

\paragraph{Objective System.} The most impactful addition was an explicit objective system that separated planning from action execution. Rather than relying on the LLM to implicitly maintain a sense of progress, the harness maintained a sequence of objectives that could be generated dynamically by the agent or pre-defined from a walkthrough source~\citep{bulbapedia2025walkthrough}. By isolating planning into a distinct phase, the agent could develop structured plans and execute them systematically, reducing the frequency of confused or aimless behavior.

\paragraph{Knowledge Base.} A persistent knowledge base---a precursor to the later \texttt{memory.json} abstraction---allowed the agent to store discoveries (locations visited, items found, NPCs encountered, strategic insights) in a flat list with importance-weighted retrieval. This gave the agent a mechanism for localizing its progress within the game narrative when generating new plans, substantially improving consistency in execution.

\paragraph{A* Pathfinding.} Integration of A* pathfinding over porymap-derived map data transformed navigation from a pure LLM reasoning task into an algorithmic one. The agent could now request navigation to specific coordinates or warp points, with the pathfinder computing optimal routes that respected tile walkability, collision boundaries, and (partially) one-way ledges. This reduced navigation failures dramatically, though limitations remained: direction-restricted tiles, dynamic NPC positions, and large maps with many warps continued to pose challenges.

\begin{figure}[h]
\centering
\includegraphics[width=\linewidth]{ch4/milestone_confidence.pdf}
\caption{\textbf{Speedrunning Harness vs.\ Model Comparisons.} Cumulative wall-clock time, actions, tokens, and cost at each milestone for five frontier models on the \pokeagent{} harness (mean $\pm$ min/max range across runs). Gemini 3 Flash completes the route fastest ($\sim$2:24 mean) but requires more actions than Gemini 3 Pro. Claude Sonnet 4.5 completes all milestones but with the highest variance and 3--4$\times$ the cost of Gemini variants. GPT-5.2 falls between the two families.}
\label{fig:ch4_harness_vs_model}
\end{figure}

The Evelord harness was used to produce the speedrunning baseline results reported in the \pokeagent{} Challenge Retrospective~\citep{karten2025pokeagent}. Figure~\ref{fig:ch4_harness_vs_model} compares five frontier models on the same harness. Several patterns emerge. Gemini 3 Flash achieves the fastest mean completion ($\sim$2:24), while Gemini 3 Pro is more step-efficient (fewer total actions). Claude Sonnet 4.5 completes all milestones but exhibits the highest variance (6:25--20:45 across runs) and 3--4$\times$ the cost of Gemini variants. GPT-5.2 falls between the two model families in both time and cost.

A consistent finding across models is that Gemini outperforms in visual reasoning, despite Claude and GPT models being generally regarded as stronger on text-only benchmarks. Gemini models exhibited more reliable visual grounding: they would almost always select Mudkip as their starter, with reasoning traces explicitly noting its type advantage against the Rock-type first gym. Claude and GPT models were less consistent in this regard, affecting downstream battle performance. While not directly relevant to our harness-level analysis---which does not involve training or weight updates---this observation nonetheless suggests that domain-specific visual knowledge is a meaningful differentiator even among frontier models.

The retrospective also evaluated three CLI agents---Claude Code, Codex CLI, and Gemini CLI---on the standardized \pokeagent{} interface. These coding-focused agents bring powerful native abstractions: context windows on the order of 200,000 tokens before compaction, strong Pok\'{e}mon world knowledge, and proprietary problem-solving scaffolding that enables hypothesis exploration and online research. Several qualitative observations stood out. Their native planning systems, while effective for software engineering tasks, operated at a level of granularity that was often insufficiently specific for RPG navigation---plans like ``go to Rustboro City'' lack the operational detail needed for reliable execution. Their proprietary scaffolding allowed them to explore hypotheses directly in the embodied setting and use online research tools to redirect when stuck, but at the cost of significantly higher token burn. These results underscored that while powerful general-purpose harnesses can make meaningful progress out of the box, the long-horizon, non-linear exploration characteristic of RPG play benefits from domain-tailored abstractions.

% ---------------------------------------------------------------------------
\subsection{Towards Subagents}\label{sec:towards_subagents}

The results from both the competition and the benchmark paper crystallized a set of architectural gaps in the Evelord harness that could not be addressed by incremental improvements. Two external trends reinforced the direction of development.

\paragraph{Model Capability Evolution.} The transition from Gemini 2.5 to Gemini 3 models brought substantial improvements in multimodal reasoning~\citep{google2025gemini3}. Actions became more visually grounded, with the agent demonstrating finer control over movement and more reliable interpretation of in-game text and menus. Tasks that had been prohibitively difficult for Gemini 2.5---such as accurately setting the bedroom clock or reading small menu text---became reliable with Gemini 3. This progress suggested that investing in richer harness abstractions was worthwhile: as the base model improved, more sophisticated tool use and planning could be leveraged effectively.

\paragraph{CLI Harness Evolution.} In parallel, the rapid maturation of CLI agent frameworks in the software engineering domain---Claude Code, Cursor, Codex CLI, Gemini CLI---demonstrated the power of subagent architectures for managing context in long-horizon tasks~\citep{cursor2026scalingagents}. These systems introduced planning tools, persistent memory, and eventually subagent delegation as first-class primitives. Their success in coding tasks provided a compelling template for embodied agent design.

\paragraph{Persistent Failure Cases.} Despite the Evelord improvements, several classes of failure persisted and pointed toward specific missing capabilities:

\begin{itemize}[nosep]
    \item \textbf{Inability to self-critique}: when the agent became stuck---repeatedly attempting the same failed navigation or battle strategy---there was no mechanism for automatic reflection. Recovery required manual intervention, such as telling the agent to try a different warp or abandon an unproductive strategy.
    \item \textbf{Objective sensitivity}: the agent was sensitive to the naming and granularity of objectives. An objective like ``route\_104\_south'' would cause the agent to navigate to the southern part of Route 104 rather than recognizing it as a distinct map section, leading to confused routing.
    \item \textbf{Pathfinding limitations}: A* pathfinding did not respect direction-restricted tiles (one-way ledges) or account for dynamic NPC positions. In the Rustboro Gym, moving trainers would invalidate computed paths, and the agent---rather than reverting to manual navigation using visual cues---would repeatedly invoke the pathfinder and fail. This required stochastic fixes (adding variance to paths) and manual intervention.
    \item \textbf{Context loss after teleportation}: when the agent fainted and was teleported to a Pok\'{e}mon Center, it had no mechanism for reviewing its completed objectives or retracing its route. It would often wander aimlessly, unable to localize its progress within the game narrative.
    \item \textbf{Large map navigation}: maps with many warps overwhelmed the agent's ability to select the correct destination, leading to loops of navigating to incorrect warps.
\end{itemize}

These failure cases motivated the design of five specific abstractions absent from the Evelord harness: \emph{subagent delegation} for context management (so that specialized reasoning tasks do not pollute the orchestrator's context window), \emph{persistent hierarchical memory} (so that the agent can record and retrieve discoveries across thousands of steps), \emph{a skill library} (so that effective behavioral and executable strategies can be saved and reused), \emph{sandboxed code execution} (so that the agent can prototype and deploy custom tools), and a \emph{general subagent primitive} (so that the agent---or an external evolutionary process---can define and launch new subagents at runtime). These abstractions collectively define the expert harness $\mathcal{H}_{\text{expert}}$, described next.

% ---------------------------------------------------------------------------
\subsection{System Design: Expert Harness}\label{sec:system_design}

The expert harness $\mathcal{H}_{\text{expert}}$ is a multi-agent orchestration system implemented primarily in \texttt{agents/PokeAgent.py} ($\sim$3,500 lines), with supporting modules under \texttt{agents/subagents/}, \texttt{utils/stores/}, and \texttt{server/}. It represents a qualitative departure from the competition-era and Evelord architectures: rather than a monolithic observe-plan-act loop, $\mathcal{H}_{\text{expert}}$ decomposes the agent into a central orchestrator that coordinates specialized subagents, persistent stores, and a sandboxed execution environment. This system is engineered specifically for the long-horizon setting described in Section~\ref{sec:long_horizon}. As we will see, it accepts a modest penalty in per-step efficiency---due to the overhead of context switching between orchestrator and subagents---in exchange for substantially better prevention of catastrophic forgetting and erroneous reasoning over extended play.

The same codebase supports three scaffold configurations, controlled by a single parameter: \texttt{"pokeagent"} ($\mathcal{H}_{\text{expert}}$) provides the full tool suite including built-in subagents, A* pathfinding, walkthrough access, and wiki lookup; \texttt{"simple"} ($\mathcal{H}_{\min}$) strips away all built-in subagents, navigation tools, and research tools, leaving only raw button presses and the persistent stores; \texttt{"autoevolve"} ($\mathcal{H}_{\text{auto}}$) extends $\mathcal{H}_{\min}$ with a \texttt{HarnessEvolver} that mutates skills, subagents, and prompts over time. The latter two configurations are central to the experimental design in Chapter~\ref{sec:autoevolve}.

\begin{figure}[t]
\centering
\includegraphics[width=\linewidth]{ch4/pokeagent_architecture.pdf}
\caption{\textbf{$\mathcal{H}_{\text{expert}}$ Multi-Agent Architecture.} The orchestrator coordinates subagents and MCP tools to play the game. Data flows left-to-right from the GBA emulator (game frames, parsed state) through the orchestrator (context management, subagent dispatch, stuck detection), which invokes either MCP tools (pathfinding, button inputs, knowledge retrieval) for direct game control or specialized subagents (battler, reflection, verification, gym puzzle, planner) for complex reasoning. A feedback loop returns compacted subagent outputs and tool results to inform subsequent decisions. \tersoo{Refine this draft with the full expert harness details (trajectory window, memory/skill/subagent stores, sandbox).}}
\label{fig:ch4_expert_architecture}
\end{figure}

\subsubsection{Architecture Overview}\label{sec:architecture_overview}

The orchestrator's main loop proceeds as follows. At each step, it (1)~fetches the current game state via the MCP \texttt{get\_game\_state} tool, (2)~assembles a structured prompt incorporating the short-term memory window, store overviews, current objectives, and the game state, (3)~sends the prompt and screenshot to the VLM, (4)~executes any tool calls from the VLM response, and (5)~updates the trajectory log and metrics. The orchestrator enforces that every step must conclude with either a \texttt{press\_buttons} or \texttt{navigate\_to} call, preventing the model from deliberating indefinitely without advancing the game.

When the orchestrator invokes a delegated subagent (e.g., the battler or planner), control transfers to a dedicated handler that runs its own multi-turn loop. Each inner turn claims a global step, fetches fresh game state, calls its own VLM with a restricted tool set, and logs to the shared trajectory. When the subagent terminates, only a compacted summary returns to the orchestrator's short-term memory---not every inner turn. This context isolation is a critical design choice: it prevents the orchestrator's context window from being polluted by hundreds of battle or planning turns, enabling it to maintain coherence over thousands of high-level decisions.

\subsubsection{Trajectory Window}\label{sec:trajectory_window}

The trajectory window is a persistent JSONL file (\texttt{trajectory\_history.jsonl}) that records every agent step as a structured entry containing the step number, action taken, pre- and post-state snapshots, reasoning trace, outcome, location, and player coordinates. This log is inspired by the recursive language model (RLM) paradigm~\citep{qu2024recursive}: rather than treating history as a flat, append-only context sequence, the trajectory is a structured, programmatically accessible workspace that subagents can query, analyze, and summarize.

The trajectory window serves two functions. First, it provides 1-step subagents---\texttt{reflect}, \texttt{verify}, and \texttt{summarize}---with the raw material for their analyses. The \texttt{process\_trajectory\_history} tool allows the VLM to run ad-hoc analysis over arbitrary step ranges, enabling the orchestrator to reason about patterns in its own behavior (e.g., detecting loops or identifying when a strategy stopped making progress). An efficient reverse-read implementation loads recent entries without scanning the full file, keeping latency low even for trajectories spanning thousands of steps.

Second, the trajectory window provides the primitive that becomes central to $\mathcal{H}_{\min}$ and $\mathcal{H}_{\text{auto}}$ in Chapter~\ref{sec:autoevolve}. In the expert harness, the trajectory is processed by built-in subagents with hand-crafted prompts. In the autoevolve setting, the same trajectory processing infrastructure is exposed to the evolutionary loop, which can create custom subagents that operate on trajectory segments in novel ways. The trajectory window thus serves as the bridge between the human-engineered and the auto-evolved harness configurations.

\subsubsection{Memory}\label{sec:memory}

The memory system implements a dual-layer architecture loosely inspired by the distinction between working memory and long-term memory in cognitive science~\citep{baddeley2000episodic, tulving1972episodic}.

\paragraph{Short-term memory.} The orchestrator maintains a \texttt{conversation\_history} buffer of the last 20 interaction turns (VLM prompts and responses). This buffer, formatted and injected into every prompt as the \textsc{Short-Term Memory} section, provides immediate context about recent actions and their outcomes. Additionally, compact tree overviews of all three persistent stores (memory, skills, subagents) are injected every step, giving the orchestrator a birds-eye view of its accumulated knowledge without consuming excessive context.

\paragraph{Long-term memory.} The persistent memory store is built on a generic \texttt{BaseStore[T]} class that provides JSON-backed CRUD operations with hierarchical path organization. Each \texttt{MemoryEntry} contains an identifier, a hierarchical path (e.g., \texttt{pokemon/team}, \texttt{events/gym\_battles}), a title, free-text content, the location and coordinates where it was created, importance (1--5), and tags. The tree overview renders entries as an indented hierarchy with \texttt{[id] title} leaves, providing a compact index that the VLM can scan to decide which entries to read, update, or delete.

The VLM interacts with memory through a unified \texttt{process\_memory} tool supporting add, read, update, and delete operations, each requiring a mandatory \texttt{reasoning} parameter that explains why the operation is being performed. Mutation tracking records every update with timestamps and old/new values, enabling post-hoc analysis of how the agent's knowledge evolves over a run. The entire store persists to \texttt{memory.json} after every write, surviving crashes and enabling cross-session continuity.

The cognitive analogy is deliberate: reading from the short-term buffer is fast and automatic (akin to working memory), while querying the long-term store requires an explicit tool call at higher latency (akin to effortful retrieval from long-term memory). This asymmetry incentivizes the agent to maintain a compact, frequently consulted overview while relegating detailed knowledge to persistent storage.

\subsubsection{Skills}\label{sec:skills}

The skill library extends the same \texttt{BaseStore} infrastructure to behavioral and executable strategies. Each \texttt{SkillEntry} contains a name, description, hierarchical path, effectiveness rating (low/medium/high), and optionally a \texttt{code} field containing executable Python. The tree overview annotates executable skills with \texttt{(run with: run\_skill)}, signaling to the VLM that these skills can be invoked directly.

Behavioral skills are textual: they encode strategies (e.g., ``when facing Rock-type gym leaders, lead with a Water or Grass type and use super-effective moves'') that the VLM can consult during reasoning. Executable skills contain Python code that runs in the sandbox (Section~\ref{sec:sandbox}) with access to game tools, enabling the agent to automate repetitive patterns---such as a custom grinding routine or a specialized battle opener---without consuming VLM inference on each step.

The VLM manages skills through \texttt{process\_skill} (CRUD) and \texttt{run\_skill} (execution). In $\mathcal{H}_{\text{expert}}$, the skill library is typically populated by the agent during gameplay as it identifies recurring patterns. In $\mathcal{H}_{\text{auto}}$, the AutoEvolve loop can also write skills, enabling the evolutionary process to encode discovered strategies as reusable code.

\subsubsection{Subagents}\label{sec:subagents}

The subagent system is the most architecturally complex abstraction in $\mathcal{H}_{\text{expert}}$. It distinguishes between two types of subagents: \emph{1-step} subagents that perform a single VLM pass without tool access, and \emph{delegated-loop} subagents that run autonomous multi-turn loops with their own restricted tool sets.

\paragraph{Built-in 1-step subagents.} Four built-in subagents operate on the current state and recent trajectory without taking game actions:
\begin{itemize}[nosep]
    \item \texttt{reflect}: diagnoses stuck or looping states by analyzing the trajectory, returning an assessment with specific recommendations for strategy changes.
    \item \texttt{verify}: independently checks whether the current objective is truly complete, returning structured evidence for and against completion with a confidence score. This prevents the orchestrator from prematurely advancing when the agent incorrectly believes a task is finished.
    \item \texttt{summarize}: produces an unbiased trajectory handoff summary, used during context compaction to preserve essential information while discarding redundant detail.
    \item \texttt{gym\_puzzle}: provides domain-specific hints for gym puzzles (e.g., Mauville's electric barriers) combined with current-state analysis.
\end{itemize}

\paragraph{Built-in delegated-loop subagents.} Two subagents run autonomous multi-turn loops:
\begin{itemize}[nosep]
    \item \texttt{battler}: handles all combat decisions---move selection, switching, item usage---for up to 200 turns with a restricted tool set (\texttt{press\_buttons}, memory/skill operations, Pok\'{e}mon lookup). The orchestrator delegates control when a battle is detected and receives a compacted summary upon completion.
    \item \texttt{planner}: manages objective creation and revision for up to 25 turns, with access to research tools, walkthrough data, and nested 1-step subagents (\texttt{summarize}, \texttt{verify}, \texttt{reflect}). The planner sees the full objective sequence across all categories (story, battling, dynamics) and can create, modify, or delete objectives.
\end{itemize}

\paragraph{Custom subagent primitive.} Beyond the built-in subagents, the VLM can create and launch arbitrary subagents at runtime via two mechanisms: (1)~registering a new subagent in the \texttt{SubagentStore} with a specified handler type, maximum turns, available tools, system instructions, and return condition, then launching it by ID; or (2)~launching an inline subagent with an ephemeral configuration. Custom subagents run through the same executor infrastructure as built-ins, with a \texttt{return\_to\_orchestrator} tool injected automatically for termination. Nesting is forbidden (custom subagents cannot launch other custom subagents), and looping subagents without available tools are rejected.

This general subagent primitive is what makes $\mathcal{H}_{\text{auto}}$ possible: the AutoEvolve loop (Chapter~\ref{sec:autoevolve}) can define new subagents with novel system instructions, tool sets, and termination conditions, effectively extending the agent's behavioral repertoire without human engineering.

\subsubsection{Sandbox and Code Execution}\label{sec:sandbox}

Two tools provide sandboxed Python execution:

\paragraph{\texttt{run\_skill}} loads a skill's \texttt{code} field from the skill store and executes it in a sandbox with restricted built-in functions (no \texttt{open}, \texttt{os}, \texttt{sys}, \texttt{subprocess}) and access to standard-library modules (\texttt{random}, \texttt{collections}, \texttt{math}, \texttt{json}, \texttt{re}, \texttt{heapq}, \texttt{itertools}, \texttt{functools}, \texttt{numpy}). Game interaction is mediated through a \texttt{tools} dictionary providing callable wrappers for \texttt{press\_buttons}, \texttt{get\_game\_state}, \texttt{get\_map\_data}, \texttt{complete\_direct\_objective}, and \texttt{process\_memory}. The skill's return value is passed back to the VLM as the tool result.

\paragraph{\texttt{run\_code}} provides the same sandbox but with read-only tools only (\texttt{get\_game\_state}, \texttt{get\_map\_data})---no \texttt{press\_buttons}. It captures \texttt{print()} output to a string buffer, returning both the printed output and any result value. This tool is designed for prototyping and debugging: the agent can test code logic before committing it to the skill store as a reusable executable skill.

The sandbox is process-local (not containerized), with security enforced through restricted built-ins. The \texttt{\_\_import\_\_} function remains accessible, allowing the agent to import additional standard-library modules as needed, but game-action tools are explicitly gated to prevent uncontrolled execution.

\vspace{1em}
Taken together, these six abstractions---trajectory window, memory, skills, subagents, sandbox, and the orchestrator that coordinates them---define a harness that is substantially more capable than the competition-era agent. The cost is complexity: $\mathcal{H}_{\text{expert}}$ comprises roughly 4,300 files (compared to 120 in the competition era) and requires significant human engineering to configure effectively. Whether this engineering effort can be partially or fully automated is the central question of Chapter~\ref{sec:autoevolve}.
